\section{Record: the circle-slice arity floor in numbers}
\label{sec:record-circle-slice}

This part is long-record only.  Nothing in the short note depends on it, and
none of it is a certified upper bound.  It records what the arity floor behind
Theorem~\ref{res:low-critical-thirteen-twentyfifths} is worth, so that a reader
deciding whether to attack the complementary regime $13/25<\mu<1$ can see the
size of the remaining slack.

At $a=1$ the constants of the failure inequalities are
$\delta=0.4586751$, $D=5.3770727$, $d_{\mathrm{low}}=2.1700771$, and
$\cosh(D/2)=e^{2}$.  Four columns are worth comparing at that area: the
earlier corpus profile $\Lambda(k,1)=\max(\delta/2+(k-1)\lambda(g),k\tau/2)$
with $\tau=0.2723415$; the hyperbolic packing floor $\Lambda_{\mathrm{pack}}$
that the previous threshold $2/5$ was built on; the linear-programming value
$M_{\mathrm{circ}}$ of the relaxation of
Lemma~\ref{res:circle-slice-packing}; and the best configuration a floating
optimiser could realise.

\begin{center}
\begin{tabular}{@{}rrrrr@{}}
\toprule
$k$ & corpus $\Lambda$ & packing & $M_{\mathrm{circ}}$ & realisable \\
\midrule
2 & 0.3103 & 0.5899 & 0.3102 & 0.3103 \\
3 & 0.4085 & 0.6878 & 0.3893 & 0.3784 \\
4 & 0.5447 & 0.7475 & 0.4450 & 0.4306 \\
5 & 0.6809 & 0.7906 & 0.4846 & 0.4694 \\
6 & 0.8170 & 0.8244 & 0.5155 & 0.4986 \\
7 & 0.9532 & 0.8521 & 0.5408 & 0.5208 \\
8 & 1.0894 & 0.8757 & 0.5622 & 0.5413 \\
10 & 1.3617 & 0.9142 & 0.5972 & 0.5759 \\
14 & 1.9064 & 0.9710 & 0.6487 & 0.6228 \\
20 & 2.7234 & 1.0298 & 0.7022 & 0.6717 \\
30 & 4.0851 & 1.0955 & 0.7624 & 0.7283 \\
100 & 13.617 & 1.2871 & 0.9428 & --- \\
\bottomrule
\end{tabular}
\end{center}

Both right-hand columns are floating searches, and the $M_{\mathrm{circ}}$
column is additionally a grid value on a $2600$-point $d$-grid and a
$1400$-point $r$-grid, so it is a lower bound for the relaxation.  The $k=2$
row shows the size of that artifact, $0.31018$ against the exactly known
$0.310338$.  Their proximity suggests a small gap and gives no rigorous
bracket.

The certified floors themselves, as integers after rounding up, read as
follows at $a=1$.

\begin{center}
\begin{tabular}{@{}lrrrrrrrr@{}}
\toprule
$x$ & 0.40 & 0.45 & 0.50 & 0.55 & 0.60 & 0.63 & 0.65 & 0.70 \\
\midrule
corpus profile & 3 & 4 & 4 & 5 & 5 & 5 & 5 & 6 \\
packing floor & 2 & 2 & 2 & 2 & 2 & 3 & 3 & 4 \\
circle-slice floor & 3 & 4 & 6 & 8 & 10 & 12 & 14 & 18 \\
\bottomrule
\end{tabular}
\end{center}

In one explicit-Euler floating replica with step $10^{-3}$ and a geometric grid
of $40$ initial areas from $10^{-6}$ to $1$, the supremum of the hitting time
falls from $0.89703$ with the packing floor to $0.65503$ with the certified
circle-slice pairs, and to $0.63003$ if the relaxation's own linear-programming
value replaces the certified pairs.  Shrinking the relaxation by the worst
measured slack factor $1.045$ moves it only to $0.60703$.  These computations
suggest a mechanism ceiling near $e^{-0.607}=0.545$; they prove no such
ceiling, and they exclude neither a stronger arity bound nor a larger certified
regime.

Two smaller leaks in the present certificate are engineering.  The dual is
restricted to $p$ radii, ten in quick mode and fourteen in full, which at
$a=1$, $k=8$ costs about one per cent against the unrestricted family; the
$a$-grid is coarse and rounded up, at a comparable cost.

One sub-statement is genuinely open and is smaller than the parent problem.
Determine $M_{\mathrm{circ}}(k,a)$ exactly, that is, maximise
$\sum_j\lambda(d_j)$ subject to $d_j\ge d_{\mathrm{low}}$ and
$\sum_jw(d_j,r)\le\pi$ for every $r>0$.  The optimiser's solution is a
spread-out density in $d$, and it appears to grow like
$(\pi/\sinh(D/2))\,c\log k$.  A closed form would determine the relaxation and
clarify the geometry of the failure configuration.  How much it would improve
the final threshold is unproved.
